\documentclass[10pt,a4paper]{amsart}
\usepackage[margin=19mm]{geometry}
\usepackage{amsmath,amssymb,mathtools}
\usepackage[scr=rsfs,cal=euler]{mathalfa}
\usepackage{xfrac}
\usepackage[unicode,hidelinks,bookmarks=false]{hyperref}
\newcommand{\ie}{i.e.\ }
\newcommand{\cf}{cf.\ }
\newcommand{\resp}{respectively\ }
\newcommand{\Subsection}{\subsection}
\providecommand{\nobreakdash}{\nobreak\hbox{-}}
\setlength{\emergencystretch}{2em}
\multlinegap0pt
\usepackage{bm}
\usepackage{bbm}
\usepackage{dsfont}
\usepackage{xspace}
\usepackage{cancel}
\usepackage{mathtools}
\usepackage{amsthm}
\makeatletter
\@ifundefined{theorem}{\newtheorem{theorem}{Theorem}}{}
\@ifundefined{lemma}{\newtheorem{lemma}[theorem]{Lemma}}{}
\makeatother
\def\dA{{\mathcal{A}}}
\def\dB{{\mathcal{B}}}
\def\dF{{\mathcal{F}}}
\def\dist{\mathop\mathrm{dist}} 						%distance
\newcommand{\Child}{\mathsf{Ch}}
\newcommand{\DD}{\mathcal{D}}
\newcommand{\Nei}{\mathsf{Nbd}}
\newcommand{\diam}{\mathrm{\, diam}}
\title[Analytic capacity and dimension of sets with plenty of big p]{Analytic capacity and dimension of sets with plenty of big projections}
\author{Damian D\k{a}browski, Michele Villa}
\date{Source: arXiv:2204.05804v2 (CC BY 4.0)}
\begin{document}
\maketitle

\noindent\emph{Source and licence.} Analytic capacity and dimension of sets with plenty of big projections. Version arXiv:2204.05804v2 (CC BY 4.0), distributed under the Creative Commons Attribution 4.0 International licence, as recorded in the arXiv licence metadata for this exact version and shown on the abstract page https://arxiv.org/abs/2204.05804v2. Full rights notice: licenses/arxiv-2204.05804-CC-BY-4.0.txt.\par\smallskip

\noindent\textit{Editorial scope.} Counted unit: the Lemma whose source label is \texttt{lem:massstays} in the article (source0.tex lines 2512--2521, source label \texttt{lem:massstays}), delivered together with the article's own complete original proof of it (source0.tex lines 2522--2587, one \texttt{proof} environment of 66 lines). No mathematics is written, repaired or re-derived: statement and proof are the article's own text line for line. Required context retained uncounted: eq:etachildrencons (lines 2304--2306). Every definition, display and label of those lines is retained unchanged. Omitted from this package and not redistributed: 1--2303, 2307--2511, 2588--3103. Nothing inside a retained excerpt is omitted or altered: every retained body is delivered exactly as the article's lines are written, so no omission receipt of any kind accompanies this package. Citations: the retained excerpts carry no citation command at all, so no bibliography entry of the article is transcribed and no reference is redistributed. Reference adaptation: none was needed and none was made; every reference inside the delivered unit resolves inside it, so no unresolved reference or citation is produced. Printed slips kept literal: no printed word, symbol, hypothesis or number of the counted statement or of its proof was altered. Layout and class: the article's own class is \texttt{amsart} with packages including inputenc, fontenc, lmodern, geometry, amsmath, amssymb, amsfonts, mathrsfs, amsthm, commath, mathtools, hyperref, doi, bigints; the wrapper is the lane's pinned \texttt{amsart} editorial wrapper at 10pt on A4, which restates the theorem-like environments the retained text needs and adds the article's own macro definitions that the retained text needs (\texttt{\textbackslash Child}, \texttt{\textbackslash DD}, \texttt{\textbackslash Nei}, \texttt{\textbackslash dA}, \texttt{\textbackslash dB}, \texttt{\textbackslash dF}, \texttt{\textbackslash diam}, \texttt{\textbackslash dist}), with extra wrapper packages (bm, bbm, dsfont, xspace, cancel, mathtools). The delivered numbering is the unit's own. Assets: no retained span contains a figure environment, a graphics-inclusion command, a PDF-page inclusion, a table or a TikZ picture; no image file is copied into this package, the package declares no asset dependency and no image file is redistributed with it. Licence: version arXiv:2204.05804v2 of this article is distributed under the Creative Commons Attribution 4.0 International licence, as recorded in the arXiv licence metadata for this exact version and shown on the abstract page https://arxiv.org/abs/2204.05804v2. A full rights notice is delivered at licenses/arxiv-2204.05804-CC-BY-4.0.txt. Characters: every retained excerpt is pure ASCII, so the delivered engine, XeLaTeX with its default Latin Modern text font, reports no missing character.
\begin{equation}\label{eq:etachildrencons}
\sum_{Q\in\Child(R)} \eta_{k}(B(Q)) = \mu_{k-1}(B(R)).
\end{equation}

\begin{lemma}\label{lem:massstays}
	If $j\ge k\ge 0$ and $Q\in\DD_k$, then
	\begin{equation}\label{eq:massstays}
	\mu_j(2B_Q)\ge \mu_k(B(Q))
	\end{equation}
	and
	\begin{equation}\label{eq:massstays2}
	\mu_j(B_Q)\le \mu_k(10B_Q.)
	\end{equation}
\end{lemma}

\begin{proof}
	In the definition of $\mu_{k+1}$ we transfer the mass of $\eta_{k+1}$ only between neighbors. It follows easily that for any family $\dF\subset\DD_k$
	\begin{equation}\label{eq:AF}
	\sum_{Q\in\dF}\mu_{k}(B(Q)) \overset{\eqref{eq:etachildrencons}}{=} \sum_{Q\in\dF}\sum_{P\in\Child(Q)}\eta_{k+1}(B(P))\le \sum_{R\in\mathcal{A}(\dF)}\mu_{k+1}(B(R)),
	\end{equation}
	where
	\begin{equation*}
	\mathcal{A}(\dF)\coloneqq \bigcup_{P\in\mathcal{F}}\bigcup_{R\in\Child(P)}\Nei(R)\subset\DD_{k+1}.
	\end{equation*}
	
	Let $Q\in\DD_k$. Set $\mathcal{A}_0(Q)\coloneqq\{Q\}$, and then inductively
	\begin{equation*}
	\mathcal{A}_i(Q) \coloneqq \mathcal{A}(\mathcal{A}_{i-1}(Q))= \bigcup_{P\in\mathcal{A}_{i-1}(Q)}\bigcup_{R\in\Child(P)}\Nei(R)\subset \DD_{k+i}.
	\end{equation*}
	Applying \eqref{eq:AF} $i$-times yields
	\begin{equation}\label{eq:AF2}
	\mu_k(B(Q))\le \sum_{P\in\dA_i(Q)}\mu_{k+i}(B(P)).
	\end{equation}
	
	Observe that if $R\in\DD$, then $\bigcup_{R'\in \Nei(R)}2B_{R'} \subset 5B_R$. It follows that for all $i\ge 0$
	\begin{equation*}
	\bigcup_{P\in\dA_i(Q)}2B_P \subset 2B_Q.
	\end{equation*}
	Indeed, this is clear for $i=0$, and then by induction
	\begin{multline*}
	\bigcup_{P\in\dA_i(Q)}2B_P = \bigcup_{P\in\mathcal{A}_{i-1}(Q)}\bigcup_{R\in\Child(P)}\bigcup_{R'\in \Nei(R)}2B_{R'} \subset \bigcup_{P\in\mathcal{A}_{i-1}(Q)}\bigcup_{R\in\Child(P)} 5B_R\\
	\subset \bigcup_{P\in\mathcal{A}_{i-1}(Q)} 2B_P \subset 2B_Q.
	\end{multline*}
	Together with \eqref{eq:AF2} this gives \eqref{eq:massstays}.
	
	Similarly, if for $\dF\subset\DD_j$ we define
	\begin{equation*}
	\mathcal{B}(\dF)\coloneqq \bigcup_{P\in\mathcal{F}}\Nei(P^1)\subset\DD_{j-1},
	\end{equation*}
	then by the definition of $\mu_j$
	\begin{multline}\label{eq:blablab}
	\sum_{Q\in\dF}\mu_j(B(Q))\le \sum_{P\in \bigcup_{Q\in\dF}\Nei(Q)} \eta_j(B(P)) \le \sum_{R\in\dB(\dF)} \sum_{P\in \Child(R)} \eta_j(B(P))\\
	 \overset{\eqref{eq:etachildrencons}}{=} \sum_{R\in\dB(\dF)} \mu_{j-1}(B(R)).
	\end{multline}
	Let $Q\in\DD_k$ and fix $j\ge k$. Set $\dB_0(Q) = \{ P\in\DD_j\ :\ B(P)\cap B_Q\neq\varnothing\}$, and for $0< i \le j-k$ we define inductively
	\begin{equation*}
	\dB_i(Q) \coloneqq \dB(\dB_{i-1}(Q)) = \bigcup_{P\in\dB_{i-1}(Q)}\Nei(P^1)\subset\DD_{j-i}.
	\end{equation*}
	By \eqref{eq:blablab}
	\begin{equation}\label{eq:blablabla}
	\mu_j(B_Q)\le \sum_{P\in\dB_0(Q)}\mu_j(B(P)) \le \sum_{P\in \dB_{j-k}(Q)}\mu_k(B(P)).
	\end{equation}
	Now, we claim that
	\begin{equation*}
	\bigcup_{P\in\dB_{j-k}(Q)}B(P) \subset 10B_Q.
	\end{equation*}
	To see this, note that each $P\in\dB_i(Q)$ has a neighbor $P'$ such that $\Child(P')\cap\dB_{i-1}(Q)\neq\varnothing.$ Hence, 
	\begin{equation*}
	\dist(P,\bigcup_{R\in \dB_{i-1}(Q)}R)\le \dist(P,P')+\diam(P')\le 3\ell(P)=15\rho^{j-i}.
	\end{equation*}
	Fix $P\in\dB_{j-k}(Q)$, and let $P=P_{j-k}, P_{j-k-1},\dots, P_{0}$ be such that for each $i$ we have $P_i\in\dB_i(Q)$ and $\dist(P_i,\bigcup_{R\in \dB_{i-1}(Q)}R)=\dist(P_i,P_{i-1})$. Then,
	\begin{multline*}
		\dist(P,\bigcup_{R\in \dB_{0}(Q)}R)\le \sum_{i=j-k}^{1} \dist(P_i,\bigcup_{R\in \dB_{i-1}(Q)}R) + \ell(P_{i-1})\\
		\le \sum_{i=j-k}^{1} 4\ell(P_{i})=\sum_{i=j-k}^{1} 20\rho^{j-i} \le 25\rho^{k} = 5\ell(Q).
	\end{multline*}
	It is easy to see that $\bigcup_{R\in \dB_{0}(Q)}R\subset 3B_Q$, and so it follows that
	\begin{equation*}
		B(P)\subset 10 B_Q.
	\end{equation*}
	Together with \eqref{eq:blablabla}, this gives \eqref{eq:massstays2}.
\end{proof}

\end{document}
